\documentclass[11pt]{article}
\input{preamble}
\renewcommand{\KomorebiKind}{Visual guide}
\renewcommand{\KomorebiSubtitle}{A visual guide to coherent realization}
\renewcommand{\KomorebiEdition}{Publication edition / October 2026}
\hypersetup{pdftitle={Komorebi Space: A Visual Guide to Coherent Realization},pdfsubject={Visual guide to coherent realizations, obstructions, and descent}}
\newcommand{\visualpage}[3]{\clearpage\phantomsection\addcontentsline{toc}{section}{#1}\section*{#1}\noindent #2\par\vspace{10pt}\begin{center}\resizebox{.98\linewidth}{!}{\input{figures/#3}}\end{center}\vspace{7pt}}
\begin{document}
\KomorebiCover
\pdfbookmark[0]{Contents}{visual-contents}
\begingroup
\makeatletter
\renewcommand{\@pnumwidth}{3em}
\renewcommand{\@tocrmarg}{4em}
\makeatother
\tableofcontents
\endgroup
\clearpage
\phantomsection\addcontentsline{toc}{section}{Reading the pictures}
\section*{Reading the pictures}
The diagrams in this guide describe mathematical objects with declared rules. A node is not a metaphor for unspecified information. An arrow has a domain, a codomain, and a mathematical interpretation. A loop can impose a new equation. A projection may preserve one quantity while changing another.

The forty plates follow four questions and their boundary cases.
\begin{enumerate}
\item Which collections of local data form a coherent realization?
\item What prevents prescribed data from extending, and what remains free after extension?
\item How do symmetries and overlap identifications survive gluing?
\item When does a small residual guarantee a nearby exact answer?
\end{enumerate}

The accompanying monograph supplies the definitions and proofs. The plates can be read in sequence or approached by topic through the contents. A diagram does not replace the hypotheses of the result it illustrates.

\begin{insight}
\textbf{The organizing distinction}\par
A local appearance is a coordinate. A realization is a compatible family of coordinates. A symmetry compares realizations. An obstruction records a failure that the permitted corrections cannot remove.
\end{insight}

\subsection*{The meaning of color and line}
In the schematic diagrams, green marks retained structure and ochre marks a discrepancy or correction. The numerical charts have their own labelled legends. Their meaning is encoded by labels, markers, and line patterns, not by color alone. Open and filled squares in Plate 16 distinguish recovered coordinates from possible support. Every chart evaluates the formula printed beside it. The three-dimensional projection plate uses genuine spatial coordinates, not a perspective sketch.

\subsection*{Mathematical scope}
The general theory of protected repair, bounded realization, correction budgets, regularization, vanishing-noise limits, solver certificates, and changes of basis is supported by Lean proofs. The Lean companion identifies the precise formal statements and explains how to reproduce their verification.

The explicit infinite diagonal examples have printed proofs, including their tail estimates, asymptotic constants, and rounded noise schedule. Those instances are not themselves formalized in Lean. The plotted enclosures illustrate the stated estimates. They do not replace the infinite arguments.

\visualpage{01 / A realization has several appearances}{A global point is an entire compatible family. Evaluating the point at a view reveals one appearance, but an individual appearance need not determine the point.}{01_views.tex}
\[
 \Gamma(F)=\{(x_v)_v:f^*x_w=x_v\text{ for every }f:v\to w\}.
\]
The arrows between fibers run opposite to the arrows in the view category. This contravariance expresses restriction: a more extensive description is mapped to the information retained by another view. Compatibility requires both direct and composite restrictions to agree.

The evaluations jointly determine the point. More strongly, a compatible family of maps from any set \(T\) factors uniquely through \(\Gamma(F)\). This is the limit universal property.
\begin{scopebox}The system of views is input. The construction does not select one universal view system, and it does not guarantee that the resulting set is nonempty.\end{scopebox}

\visualpage{02 / An obstruction is a quotient of actual defects}{Start with a partial section. Extend it arbitrarily as raw data, measure its defect, and ask whether an admissible correction can cancel that defect.}{02_extension.tex}
\[
 \mathfrak o_U(s)=[\delta\widetilde s]\in R/\delta(\ker p_U),
 \qquad
 \mathfrak o_U(s)=0\iff\Ext_U(s)\ne\varnothing.
\]
The denominator is the set of residuals produced by changes invisible on the prescribed views. Two raw lifts of the same partial section differ by such a change. Their obstruction classes therefore agree.

Once an extension exists, the invisible coherent changes form the ambiguity group \(H_U=\ker\delta\cap\ker p_U\). Existence is controlled by the quotient. Uniqueness is controlled by this different kernel.
\begin{scopebox}The obstruction is defined before testing solvability. Nevertheless the quotient construction does not provide a general decision procedure for membership.\end{scopebox}

\visualpage{03 / The smallest nontrivial extension problem}{The coarse value is prescribed. The finer value must be an integer whose double is that coarse value. Filled points extend. Crosses do not.}{13_parity.tex}
\[
 x_0=2x_1,\qquad \Ext_{\{0\}}(k)=\{m\in\Z:2m=k\}.
\]
For an even seed, the extension is unique. For an odd seed, it is absent. The seed \(2\) extends as \((2,1)\), so the construction is not limited to the zero realization.

Adding an independent unprescribed integer coordinate changes the number of extensions of every even seed from one to infinitely many. It leaves the parity obstruction unchanged. This separates feasibility from ambiguity in a model that can be checked completely.
\begin{scopebox}The ambient realization space is nonempty, but not every prescribed partial datum belongs to the image of restriction.\end{scopebox}

\visualpage{04 / Exact maps commute with restriction}{An exact transport is not just a map at each level. Every adjacent square must commute. The same requirement then holds over any finite interval of the tower.}{03_tower.tex}
\[
 (\partial f)_n=\rho^Y_nf_{n+1}-f_n\rho^X_n,
 \qquad \partial f=0\iff f\text{ is exact}.
\]
Exact transports induce maps on coherent inverse-limit points. Their residual obeys a precise Leibniz identity under composition:
\[
 \partial(gf)_n=(\partial g)_nf_{n+1}+g_n(\partial f)_n.
\]
A cofinal schedule preserves the coherent point space. It need not preserve the numerical residual of an arbitrary inexact family. Subsampling can cancel adjacent defects.
\begin{scopebox}A fixed-level tower category is not automatically the full procategory. Reindexing and eventual equality of maps require their own construction.\end{scopebox}

\visualpage{05 / Repair changes the residual coordinate}{Under the split-tower hypotheses, a candidate has two exact coordinates: protected information and actual compatibility residual. Repair sets the second coordinate to zero.}{04_split.tex}
\[
 \Phi(f)=(p_Lf,\partial f),\qquad
 \Phi^{-1}(s,r)=Ps+Ar,\qquad P=I-A\partial.
\]
The chosen target sections define the innovation observable \(p_L\). It retains the seed and the new target-kernel components at every level. The correction \(Ar\) changes neither of these protected components.

A propagation gap, \(SB_X<1\), makes the recursive right inverse bounded. Its gain is at most \(S/(1-SB_X)\). This is a derived repair map, not an assumed existence statement.
\begin{scopebox}The coordinate equivalence is not generally an isometry. The protected observable and the repaired answer can change when the splitting changes.\end{scopebox}

\visualpage{06 / Retain the detail that coarsening removes}{A right inverse of the coarsener supplies a chosen lift. Subtracting the lifted coarse part leaves the detail. Repeating this construction gives an exact finite telescoping decomposition.}{05_cascade.tex}
\[
 C_jJ_j=I,\qquad h_j=(I-J_jC_j)x_j,\qquad
 x=S_Nx+Q_Nx.
\]
Here \(S_Nx\) is the sum of the retained details lifted back to the original level. The remainder \(Q_Nx\) is the coarse part also restored to the original level. Both terms live in the same space, so their difference and norm are meaningful.

Every finite identity is exact. Passing to an infinite reconstruction is a different question: the restored remainder must converge to zero in the chosen topology.
\begin{scopebox}A coarse observation alone can lose information. The pair consisting of coarse data and retained detail is lossless under the stated splitting identity.\end{scopebox}

\visualpage{07 / Every coordinate can recover while the norm does not}{For the shift cascade on bounded sequences, the remainder consists of the tail in its original coordinates. Every fixed coordinate eventually disappears from the remainder.}{14_tail.tex}
\[
 (Q_Nx)_i=\begin{cases}0&i<N,\\x_i&i\ge N.\end{cases}
 \qquad \norm{Q_N\mathbf1}_\infty=1.
\]
The displayed constant sequence is recovered at each fixed coordinate, but the uniform error never decreases. Its unrecovered mass moves farther out instead of shrinking.

In this model, norm recovery occurs precisely for \(c_0\), the sequences tending to zero. More generally, uniformly bounded remainder projections give norm recovery exactly on the closure of the finite core.
\begin{scopebox}Coordinatewise convergence, uniform convergence, and existence of a coherent inverse-limit point are different assertions. No one of these topologies should be substituted for another.\end{scopebox}

\visualpage{08 / A groupoid remembers more than an orbit set}{A cocycle records coherent transition rules. Gauge arrows compare those rules. Passing to an orbit set forgets the comparison witnesses and the automorphisms.}{06_groupoid.tex}
\[
 c_{gf}=c_ff^*(c_g),\qquad
 (a\cdot c)_f=a_vc_f[f^*(a_w)]^{-1}.
\]
For a single \(C_2\) loop with coefficients in \(S_3\), cocycles correspond to involutions. There are four involutions but only two conjugacy classes: the identity and the three transpositions. The identity object's automorphism group is \(S_3\), which is noncommutative.

A nonidentity loop cocycle is coherent transition data. Its associated regular twisted section equation \(x=px\) has no solution. Coherence of rules is not the same as existence of a global trivialization.
\begin{scopebox}A single orbit does not imply a unique comparison map or a trivial symmetry group.\end{scopebox}

\visualpage{09 / Cover the equations that must be checked}{Covering all objects is weaker than covering all arrows. For nonabelian cocycles, the composition equations involve composable pairs, which may require three view objects at once.}{07_cover.tex}
\[
 \forall(v\xrightarrow f w\xrightarrow g z),\quad
 \exists i:\ \{v,w,z\}\subseteq U_i.
\]
On three ordered views, all proper two-view patches cover every individual arrow but miss the full composition support. On four ordered views, the four proper three-view patches cover every composable pair. No patch has to be the whole space.

This composition-cover rule generates a Grothendieck topology on full patches. The actual cocycle pseudofunctor satisfies mathlib's standard stack property for that topology. Objects glue up to the supplied overlap isomorphisms, and compatible morphisms glue as well.
\begin{scopebox}The stack theorem is for this specified coefficient system and site. It is not an assertion that every possible Komorebi doctrine has effective descent.\end{scopebox}

\visualpage{10 / Keep the identification of the boundary}{A realization over a boundary is an object together with an actual identification of its image with that boundary. A comparison is allowed only when it respects that identification.}{08_marked.tex}
\[
 \mathcal E_b(R)=\{(x,\phi):\phi:Rx\xrightarrow{\sim}b\},
 \qquad \psi R(a)=\phi.
\]
The marked realization groupoid retains automorphisms invisible at the boundary. Replacing a boundary by an isomorphic one transports the problem rather than forcing literal equality of representatives.

When a square of restriction functors commutes up to a specified natural isomorphism, equivalences on its horizontal sides induce an equivalence of the marked realization groupoids. The comparison isomorphism is used in the construction.
\begin{scopebox}Presentation invariance is relative to actual equivalences, pulled-back coefficients, and the stated domain hypotheses. Arbitrary changes of representation are not automatically lossless.\end{scopebox}

\visualpage{11 / Loops select the feasible anchor}{Chosen routes reconstruct vertex data from a root value. Every extra edge closes a loop. A reconstructed family is exact precisely when the anchor is fixed by all resulting loop operators.}{09_holonomy.tex}
\[
 K=\bigcap_{(i,j)}\ker(K_{ij}-I),\qquad
 \{\text{parallel families}\}\simeq K.
\]
The reflection on the right fixes one line. The anchor may vary on that line, but its transverse component cannot be part of an exact parallel family. Orthogonal projection to \(K\), followed by route reconstruction, gives a specific exact repair.

In a nonisometric connection, route products amplify defects. A flat scalar chain can have edge residual one and root-preserving repair error \(2^n-1\). Flatness alone is not a quantitative estimate.
\begin{scopebox}Orthogonality is defined on the anchor fiber. The reconstructed graph repair is not asserted to be a nearest point for the global maximum norm.\end{scopebox}

\visualpage{12 / Exactness and stable recovery are different}{The orthogonal projection onto a closed kernel exists. A uniform estimate of the correction size by the residual requires more: the range of the defect operator must be closed.}{10_closed_range.tex}
\[
 \ran D\text{ closed}\iff
 \exists C\ge0\ \forall x:\norm{x-P_{\ker D}x}\le C\norm{Dx}.
\]
For \(D(x_n)=(x_n/n)\) on \(\ell^2\), the unit vectors remain distance one from the zero kernel while their residuals tend to zero. The chart shows these two quantities at the indicated indices. There is no uniform inverse gain.

This failure cannot be removed merely by choosing a nonlinear rule for selecting exact answers. If any such rule had the same uniform distance estimate, best approximation would give the forbidden projection estimate.
\begin{scopebox}A zero residual characterizes exactness. A small residual gives proximity only after an appropriate stability certificate has been proved.\end{scopebox}

\visualpage{13 / Correct the overlap inside the locally exact space}{First project to the kernel of the local defect. Then solve the overlap problem within that kernel. This nested construction preserves local exactness during the second correction.}{11_descent.tex}
\[
 L=\ker D,\qquad M=\ker D\cap\ker C,
 \qquad P_MP_L=P_M.
\]
The picture uses \(D(x)=x_3\), \(C(x)=x_2\), and \(x=(2,2,3)\). Its first correction has length three. The second has length two. They are perpendicular, so the total squared error is thirteen. Both stages are nontrivial, and the joint answer is \((2,0,0)\).

Stability requires control of the local operator and of the overlap operator restricted to \(L\). With closed local range, closed range of that restricted matching map is equivalent to closed range of the joint defect.
\begin{scopebox}Projecting independently to \(\ker D\) and then to \(\ker C\) is a different procedure. The second unrestricted projection need not preserve the first constraint.\end{scopebox}

\visualpage{14 / An exact space can move without being preserved}{For \(D_t(x,y)=x+ty\), the kernel is the moving line \(x=-ty\). The left panel distinguishes orthogonal repair from a splitting-dependent transport of an old exact vector.}{12_moving_line.tex}
\[
 q_t=\frac1{1+t^2}\begin{pmatrix}t^2&-t\\-t&1\end{pmatrix},
 \qquad\norm{q_t-q_0}=\frac{|t|}{\sqrt{1+t^2}}.
\]
The right panel compares the exact norm with the general upper bound \(|t|\). The general theorem controls projector motion by \(\max(C,C')\norm{D-D'}\), where each model has its own residual certificate.

At the plotted values \(t=3/4\), \(b=2\), the transported point is \((-3/2,2)\), while the nearest point is \((-24/25,32/25)\). The projector gap is exactly \(3/5\). Both points lie on the new exact line, but they implement different choices.
\begin{scopebox}Continuity of the exact space is not literal equality of kernels. Both certificate control and the distinction between orthogonal projection and chosen transport matter.\end{scopebox}


\visualpage{15 / The preserved information changes the answer}{The raw point is \((2,0,5)\). Compatibility requires equality of the first two coordinates. The third coordinate remains five and is omitted from the drawing.}{15_protected_geometry.tex}
\[
 y_{\rm nearest}=(1,1,5),\qquad
 y_{\rm protected}=(2,2,5).
\]
Unrestricted orthogonal repair has squared error two. Preserving the first and third coordinates forces squared error four. Both answers are exact. Only the second belongs to the declared protected slice.

For a bounded observable \(p\), permissible corrections lie in \(M=\ker p\). The correct analytic operator is \(B=D|_M\). Feasible answers have a unique nearest representative, and a uniform residual-to-error bound exists precisely when \(\ran B\) is closed.
\begin{scopebox}The relevant range is \(D(\ker p)\). A small residual alone does not certify that the preserved information is correct, or even that protected repair is feasible.\end{scopebox}

\visualpage{16 / The actual dyadic remainder has sparse support}{In the dyadic tower lift, every even coordinate is zero. An odd coordinate is a section applied to a coarser value. Repeating this rule locates every coordinate at which a remainder can survive.}{16_dyadic_support.tex}
\[
 (Q_Nr)_k\ne0\ \Longrightarrow\ 2^N\mid(k+1).
\]
The filled squares show \emph{possible} support, not necessarily nonzero values. At depth two, only indices \(3,7,11,15,\ldots\) can remain. At depth three, only \(7,15,23,\ldots\) can remain. Coordinate \(k\) has recovered once \(N>\nu_2(k+1)\).

The monograph gives the parity induction. It respects the changing residual spaces at each depth. For scalar identity bonds and a constant-one residual, each displayed surviving value is \(2^N\). This also shows why fixed-coordinate recovery alone is not a uniform-norm statement.
\begin{scopebox}This is the actual dyadic support rule. Plate 07 is the different shift cascade. The two models illustrate the same topological distinction without having the same support pattern.\end{scopebox}

\visualpage{17 / Flat transport can amplify a small defect}{On the flat scalar chain with coordinates \(e_i(z)=2^iz\), choose \(x_i=8\cdot2^i-1\). Exact repair preserving the root value seven is \(y_i=7\cdot2^i\).}{17_amplification.tex}
\[
 \norm{dx}_\infty=1,\qquad
 \norm{x-y}_\infty=2^n-1.
\]
At eight edges, the measured defect is still one while the root-preserving correction has norm 255. The vertical axis is logarithmic and explicitly labelled. Route length grows only linearly.

The connection is flat: no inconsistent loop is responsible for the effect. The amplification comes from the products of nonisometric transport maps. An estimate must include their gains rather than only the number of edges.
\begin{scopebox}This is not a counterexample to the amplified repair theorem. It is a concrete reason for the transport factors in that theorem. Nor is root-preserving repair claimed to minimize the global maximum-norm error.\end{scopebox}

\visualpage{18 / A rank jump breaks projector continuity}{Both compared families have operator error \(\varepsilon\) from their own reference at zero. Their exact spaces behave differently. The axis moves toward smaller errors from left to right.}{18_rank_loss.tex}
\[
 D_t=(1\ \ t),\qquad
 A_t=\begin{pmatrix}1&0\\0&t\end{pmatrix}.
\]
The line kernel of \(D_t\) moves continuously, with gap \(|t|/\sqrt{1+t^2}\). For \(A_t\), the kernel at zero is the vertical axis, but at every nonzero \(t\) it is zero. Its projector therefore stays distance one from the reference projector.

For small nonzero \(t\), the diagonal family's residual gain is at least \(1/|t|\). The required certificate deteriorates just as the operator error shrinks. The two-sided motion theorem retains this dependence.
\begin{scopebox}A small operator perturbation is not a rank-stability theorem. Controlled residual certificates, or a proved split-surjective radius, supply the missing hypothesis.\end{scopebox}


\visualpage{19 / A later choice can create an obstruction}{An initial prescribed value admits infinitely many completions. Prescribing a middle coordinate removes that freedom, but it may also select an impossible branch.}{19_branches.tex}
\[
 ((7,0),(7,2n),(7,n)),\qquad n\in\Z.
\]
The first view is fixed at \((7,0)\). A proposed middle value \((7,m)\) extends exactly when \(m=2n\). Even values have a unique completion. Odd values have none. The initial problem itself remains solvable.

For nested prescribed domains \(U\subseteq W\), the ambiguity group shrinks and the quotient obstruction has a forgetful map \(Q_W\to Q_U\). A nonzero obstruction at the earlier stage persists, but a zero earlier obstruction need not make every later choice feasible.
\begin{scopebox}Solving the first stage is different from selecting a compatible solution at every stage. The figure shows an integer problem; division by two over another field changes its feasibility and ambiguity.\end{scopebox}

\visualpage{20 / Finite observations can be put on one schedule}{Different output coordinates may have independently chosen sufficient input depths. Their running maximum gives a monotone schedule without changing the realized transformation.}{20_schedule.tex}
\[
 \phi(n)=\max_{i\le n}k_i,\qquad
 f_n=a_n\rho_{k_n,\phi(n)}.
\]
The larger input is first restricted to the originally sufficient depth. Under surjectivity of the source's point projections, output coherence on global inputs proves the required finite-stage naturality equations.

Consequently, eventual transport classes correspond exactly to functions whose every output coordinate factors through one finite input view. That depth is uniform over all inputs for the chosen output coordinate. A function asking whether an infinite binary sequence is entirely zero does not have this property.
\begin{scopebox}The characterization is conditional on the stated source observability. Neither finite observation nor monotonicity supplies an efficiency bound on the input depth.\end{scopebox}

\visualpage{21 / Repair and representation change need not commute}{The raw candidate is \((0,1)\), with residual given by its second coordinate. The shear \(S(u,r)=(u+2r,r)\) preserves that residual but changes the chosen exact coordinate.}{21_mismatch.tex}
\[
 P_BSx-SP_Ax=(SA-BC)\,dx.
\]
Repairing after the shear yields \((2,0)\). Repairing before it yields \((0,0)\). Both are exact. A test of final residual alone therefore cannot tell the two selected answers apart.

In split coordinates a compatible representation change has the form \((u,r)\mapsto(Su+Mr,Cr)\), with \(M=SA-BC\) taking values in the target exact space. The two routes commute on all inputs exactly when the chosen residual sections commute as well.
\begin{scopebox}Compatibility transports equations correctly. It does not by itself transport a selected repair, a protected observable, or a norm without change.\end{scopebox}

\visualpage{22 / A shrinking margin makes repair expensive}{For \(D_t(a,b)=(1-t)a\) and \(0\le t<1\), the exact space stays the vertical axis, but the least nonnegative residual-to-distance gain is \(1/(1-t)\).}{22_gain_budget.tex}
\[
 \dist((a,b),\ker D_t)=|a|
 =\frac{|D_t(a,b)|}{1-t}.
\]
The old-kernel-preservation theorem gives the same denominator \(1-C\varepsilon\). The strict radius matters. At the endpoint \(t=1\), the operator is zero and the kernel suddenly becomes the whole plane.

Measurement uncertainty contributes its own additive budget. With an uncertain operator, a reference-space guarantee also contains a term proportional to the input norm. A zero measured residual for the changed operator need not imply exactness for the reference operator.
\begin{scopebox}The endpoint does not have infinite least gain. Its exact space changes. The diverging curve is the gain on the approaching family with the old kernel.\end{scopebox}

\visualpage{23 / Close the signal loop, then choose an anchor}{The raw edge data \((1,0,0,0)\) cannot be the differences around a closed four-cycle. Their sum is one. Nearest unconstrained repair subtracts the mean from every edge.}{23_cycle_primitive.tex}
\[
 \widehat r=(3/4,-1/4,-1/4,-1/4).
\]
Anchoring the first vertex at zero gives \((0,3/4,1/2,1/4)\). Subtracting its mean gives the mean-zero primitive \((-3/8,3/8,1/8,-1/8)\). The two primitives have identical edge differences. The fifth plotted point repeats the first and closes the cycle.

Preserving the first three raw edges forces the different repair \((1,0,0,-1)\). Its squared residual correction is one, versus one quarter for the unconstrained answer. These are different mathematical optimization problems, not competing numerical approximations to one answer.
\begin{scopebox}Correcting inconsistent edge data changes the realization problem. Choosing an anchor selects a primitive after the corrected edge problem is feasible.\end{scopebox}

\visualpage{24 / Nontrivial holonomy can leave an exact direction}{Two identity edges reconstruct all three vertices from an anchor. The closing edge applies a reflection. Closure permits every horizontal anchor, but excludes its vertical component.}{24_reflection_cycle.tex}
\[
 J(a,b)=(a,-b),\qquad
 \operatorname{Fix}J=\{(a,0):a\in\R\}.
\]
For raw vertices \((2,1),(4,-1),(3,2)\), projecting the original anchor and reconstructing gives \((2,0)\) at every vertex. Its squared error is eleven. Nearest repair in the declared six-dimensional Euclidean norm gives \((3,0)\) at every vertex, with squared error eight.

The exact family is the same in both constructions. The selection rule is not. An explicit Gram calculation in the monograph gives least residual-to-distance gain one for its one-orientation Euclidean residual norm.
\begin{scopebox}That gain belongs to the norm specified in this worked example. Duplicating both edge orientations or changing the product norm changes the numerical constant.\end{scopebox}

\visualpage{25 / A summable path need not have a fixed rate}{The residual decays like $1/(n+1)$. Its ratios tend to one, so no uniform contraction factor below one applies. The movements have a summable telescoping formula.}{25_summable_path.tex}
\[
 f_n=(s,(n+1)^{-1}),\qquad
 \dist(f_n,f_{n+1})=\frac1{(n+1)(n+2)}.
\]
In the maximum norm, the distance to $(s,0)$ is exactly the sum of all future movements. The residual series itself diverges, but it is not the series required by the theorem. The weighted movement series converges.

For a complete metric candidate space, lower semicontinuous nonnegative residual, vanishing residual values, and summable step bounds, an exact limit exists. The tail of the step bounds is an error certificate.
\begin{scopebox}The metric belongs to the candidate space. Small measured residuals alone do not prove that the candidates form a Cauchy sequence.\end{scopebox}

\visualpage{26 / Temporary increases do not prevent exactification}{Here every even-to-odd step doubles the residual. The next step multiplies it by one eighth. The two-step product is one quarter.}{26_residual_bursts.tex}
\[
 Q_{2k}=4^{-k},\qquad Q_{2k+1}=2\cdot4^{-k}.
\]
The movement coefficients are $C_{2k}=1$ and $C_{2k+1}=7/8$. Their weighted sum is
\[
 \sum_{n\ge0}C_nQ_n=\frac{11}{3}.
\]
The path therefore converges even though its residual increases infinitely often. A variable-rate certificate controls the prefix products and total movement, not merely whether each plotted step points down.
\begin{scopebox}A decaying product by itself is insufficient. The movement budget must also be summable. Sampling only even steps would hide the real excursions.\end{scopebox}

\visualpage{27 / The same residual history can lead to different answers}{The coordinate $s$ is fixed and not displayed. The remaining coordinates follow $u=a(1-2^{-n})$ and $r=2^{-n}$ from one common starting point.}{27_path_limits.tex}
\[
 f_n^{(a)}=(s,a(1-2^{-n}),2^{-n})\longrightarrow(s,a,0).
\]
All paths are protected and have the same exponentially decaying residual. Their total movement is finite. Their exact limits differ because $u$ is not protected.

To identify the limits of two convergent paths, a bound on their mutual distance tending to zero is sufficient. In this example their mutual distance instead tends to $|a-b|$.
\begin{scopebox}Uniqueness of the limit of one sequence is not uniqueness of the exact answer across different valid procedures.\end{scopebox}

\visualpage{28 / Smaller residuals can require escaping corrections}{The first $N$ equations are repaired exactly while the protected value stays one. The squared correction size grows, even as the infinite remaining residual tends to zero.}{28_escaping_repair.tex}
\[
 d(u,t)_n=\frac{u_n+t}{n},\qquad
 f_N=(-\mathbf1_N,1)\in\ell^2\oplus\R.
\]
An exact protected answer would require $u_n=-1$ for every $n$, outside $\ell^2$. The ordinary inverse limit of finite solutions still contains that sequence. What fails is realization in the specified Hilbert space.
\begin{scopebox}The shaded interval is a pair of proved bounds on an infinite tail. It is not an extrapolation from a computed finite truncation.\end{scopebox}

\visualpage{29 / A cost geometry is not an isomorphism test}{Arrow costs are added along paths. The diagram shows the generators of a free path category. An isolated vertex has infinite distance from the others.}{29_transport_cost.tex}
\[
 d_W(A,B)=\inf_{f:A\to B}W(f),\qquad
 d_W^{\mathrm{sym}}(A,B)=\max\{d_W(A,B),d_W(B,A)\}.
\]
The cheapest paths from $B$ to $C$ and from $C$ to $B$ cost two and five, respectively. The zero-cost arrows between $A$ and $B$ have no imposed inverse relations. Their composite is a path of length two, not an identity. Thus zero symmetric distance does not establish categorical isomorphism.

In an additive category, a cost that vanishes on every zero arrow collapses all object distances. Compatibility alone cannot supply a separated information geometry.
\begin{scopebox}The arrow cost is extra data and must satisfy the composition inequality. An arbitrary operator defect norm is not automatically such a cost.\end{scopebox}

\visualpage{30 / Associativity can coexist with an irreducible phase}{The two-bit phase satisfies every triple composition equation, but it is asymmetric in its two arguments. Every scalar coboundary on the abelian label group is symmetric.}{30_phase.tex}
\[
 c(x,y)+c(x+y,z)=c(y,z)+c(x,y+z).
\]
The matrices $X=\left(\begin{smallmatrix}0&1\\1&0\end{smallmatrix}\right)$ and $Z=\left(\begin{smallmatrix}1&0\\0&-1\end{smallmatrix}\right)$ satisfy $ZX=-XZ$. With $U_x=X^{x_1}Z^{x_2}$, their multiplier is $(-1)^{c(x,y)}$.

No scalar change removes that multiplier. Retaining an extra phase bit instead produces an eight-element group with an ordinary faithful matrix representation. The enlarged group is not the original four-label group.
\begin{scopebox}This pair-phase model and the earlier arrow-labelled cocycle stack have different data and degrees. Neither is obtained from the other by renaming its obstruction.\end{scopebox}

\visualpage{31 / Bounded finite energy produces a realization}{The same diagonal operator can have attainable and unattainable data. The minimum correction energies, not finite solvability alone, distinguish them.}{31_nested_energy.tex}
\[
 T a=(a_n/n),\qquad
 r_n=1/n^2:\ m_N^2=\sum_{n=1}^{N}n^{-2},\qquad
 w_n=1/n:\ m_N^2=N.
\]
For nested nonempty closed affine constraints in a Hilbert space, the minimum norm corrections satisfy
\[
 \norm{a_m-a_n}^2=\norm{a_m}^2-\norm{a_n}^2\qquad(m\ge n).
\]
A uniform energy bound makes them Cauchy and produces the least norm point satisfying all constraints. The escaping case lacks exactly that bound.
\begin{scopebox}The theorem concerns minimum norm answers to nested affine constraints in one common Hilbert space. Arbitrary bounded approximate sequences need not converge in norm.\end{scopebox}

\visualpage{32 / A witness certifies the optimum over a whole ball}{The correction ball maps to an ellipse in residual space. The residual from its nearest point is normal to the supporting line.}{32_budget_certificate.tex}
\[
 \beta(M)=\max_{\norm z\le1}
       \{\langle r,z\rangle-M\norm{T^*z}\}.
\]
For $T=\operatorname{diag}(1,1/2)$ and $r=(1,1)$, the correction $a=(1/2,2/5)$ has budget $M=\sqrt{41}/10$. Its residual is $v=(1/2,4/5)$ and $T^*v=a$. The normalized residual is the sharp witness, proving $\beta(M)=\sqrt{89}/10$.

The certificate excludes every better correction in the ball. It does not rely on a sampled grid of candidates.
\begin{scopebox}For protected repair, the adjoint belongs to the actual restricted operator on $\ker p$, with its declared metric.\end{scopebox}

\visualpage{33 / A larger correction amplifies model uncertainty}{The nominal system has $T=1$ and $r=1$. Its exact fixed-budget profile is $\max(1-M,0)$. Model and measurement errors widen the guaranteed band.}{33_budget_uncertainty.tex}
\[
 |\beta_{T,r}(M)-\beta_{\widetilde T,\widetilde r}(M)|
 \le\delta+\eta M.
\]
The plotted bounds use $\delta=1/20$ and $\eta=1/10$. For $M=1/2$, a nominal residual of $1/2$ gives a robust lower bound $2/5$. A true residual below $3/10$ is impossible at that budget under the stated errors.

Increasing the budget may lower the nominal optimum while increasing the error bar. A zero lower bound is not evidence that an exact repair exists.
\begin{scopebox}The compared operators act on one fixed correction space. An uncertain preservation kernel requires a separate analysis.\end{scopebox}

\visualpage{34 / Regularization limits the noise response}{The scalar noise response peaks where the singular value matches the square root of the regularization parameter.}{34_noise_filter.tex}
\[
 a_\alpha=(T^*T+\alpha I)^{-1}T^*r,\qquad
 \norm{a_\alpha(r+e)-a_\alpha(r)}\le\frac{\norm e}{2\sqrt\alpha}.
\]
A positive parameter always gives a unique protected correction. Reducing it lowers the residual but permits a larger correction and a larger worst-case noise gain. Convergence as noise vanishes requires a parameter choice balancing both effects, as well as an attainable projected exact datum.
\begin{scopebox}The noise bound is uniform at a fixed parameter. It is not evidence that an exact protected answer exists when the original residual is outside the range.\end{scopebox}

\visualpage{35 / Smoothing the correction does not remove escape}{The correction is now a smooth spectral profile rather than a hard prefix. It still tends coordinate by coordinate to the forbidden constant sequence.}{35_regularized_escape.tex}
\[
 a^{(m)}_n=\frac{m^2}{m^2+n^2},\quad
 v^{(m)}_n=\frac{n}{m^2+n^2},\quad\alpha=m^{-2}.
\]
Every individual correction belongs to $\ell^2$. Its squared energy satisfies $m/4\le A_m\le4m/3$, while the residual square satisfies $1/[4(m+1)]\le B_m\le2/m$. Thus the corrections escape even though the residual and the weighted penalty both vanish.
\begin{scopebox}A penalty term $\alpha\norm a^2$ tending to zero does not bound $\norm a$ when $\alpha$ is also tending to zero.\end{scopebox}

\visualpage{36 / The best approximation can improve without becoming realizable}{The product of correction norm and residual norm tends to a nonzero constant. The soft regularized corrections improve that constant but do not create a finite-budget exact answer.}{36_budget_product.tex}
\[
 \frac{A_m}{m}\longrightarrow\frac\pi4,\qquad
 mB_m\longrightarrow\frac\pi4,\qquad
 M\beta(M)\longrightarrow\frac\pi4.
\]
The hard-prefix constant is one. The finite plot uses each method's natural scale, not identical budgets at each marker. The proved asymptotic is the statement that compares large matched budgets.

For the distinct attainable datum $(1/n^2)$, the same residual formula becomes a vanishing regularization bias, and a deterministic noise choice yields an explicit $\delta^{1/3}$ bound.
\begin{scopebox}The plotted intervals enclose infinite quantities using exact rational partial sums and proved tails. The leading constants are proved by Riemann sums and integrals, not fitted from the graph.\end{scopebox}


\visualpage{37 / A stopping tolerance needs its parameter}{A normal-equation residual is converted into an error bound by the declared regularization parameter. The same tolerance gives different guarantees at different parameters.}{37_solver_amplification.tex}
\[
 g=T^*(Ta-r)+\alpha a,\qquad
 \norm{a-a_\alpha}\le\frac{\norm g}{\alpha},\qquad
 \norm{T(a-a_\alpha)}\le\frac{\norm g}{2\sqrt\alpha}.
\]
The curves hold the normal-residual norm at $10^{-5}$. They are universal upper bounds, not measured errors. With $T=0$, $r=0$, $a_n=1$ and $\alpha_n=n^{-2}$, the fit is zero and $g_n\to0$, yet the error stays one.
\begin{scopebox}A vanishing normal residual without a parameter-scaled bound does not certify convergence. These estimates concern the regularized optimizer, not automatic existence of an unregularized realization.\end{scopebox}

\visualpage{38 / A checkable gap brackets an unknown optimum}{A proposed correction and a dual witness certify objective accuracy without computing the minimizer first.}{38_objective_certificate.tex}
\[
 J_\alpha(a)-\mathcal D_\alpha(z)
 =\norm{Ta-r+z}^2+\frac{\norm{T^*z-\alpha a}^2}{\alpha}.
\]
On the protected line $a=(t,t)$, with $d(x,y)=x$, $p(x,y)=x-y$, $r=1$ and $\alpha=1$, the proposal $t=1/4$ has objective $11/16$. Its residual witness gives the lower bound $21/32$. The true optimum is $2/3$, between them.

The certified gap is $1/32$, while the actual suboptimality is $1/48$. A certificate need not equal the quantity it bounds.
\begin{scopebox}Every dual witness gives a valid lower bound, but a poor witness may be loose. The residual witness recovers the normal-residual certificate exactly.\end{scopebox}

\visualpage{39 / Coordinates can shrink a residual without improving an answer}{Rescaling the basis changes the normal-equation covector. Its Euclidean coordinate length is not the Hilbert error certificate.}{39_covector_metric.tex}
\[
 G=B^{\mathsf T}B,\quad q=(S^{\mathsf T}S+\alpha G)c-S^{\mathsf T}r,\quad
 g=BG^{-1}q,\quad \norm g^2=q^{\mathsf T}G^{-1}q.
\]
For the same protected proposal, replace $B=(1,1)^{\mathsf T}$ by $sB$. The raw covector magnitude is $s/4$, which can be arbitrarily small. Its Gram-corrected norm remains $1/\sqrt{32}$.
\begin{scopebox}The basis must span the entire declared correction kernel. A stationary point in a smaller subspace is not automatically the optimizer in the full kernel.\end{scopebox}

\visualpage{40 / Solver accuracy must follow the regularization schedule}{The total certified error has three terms: regularization bias, data noise, and incomplete solution of the normal equation.}{40_stopping_schedule.tex}
\[
 \norm{a-a^\dagger}\le
 \norm{a_\alpha(r)-a^\dagger}
 +\frac{\delta}{2\sqrt\alpha}+\frac{\epsilon}{\alpha}.
\]
For the attainable diagonal datum, use $m=\lceil\delta^{-2/3}\rceil$, $\alpha=m^{-2}$, and $\epsilon\le\frac14\alpha\delta^{1/3}$. The displayed three budgets all vanish. Omitting the solver term would assume an exact optimizer was already available.
\begin{scopebox}This limit needs an attained projected datum. For the escaping datum, solver accuracy only tracks the regularized path more closely. It does not make that path converge to a missing Hilbert realization.\end{scopebox}

\clearpage\phantomsection\addcontentsline{toc}{section}{The complete picture}
\section*{The complete picture}
The set-valued layer asks whether the views agree. The additive layer measures their failure to agree and restricts which corrections are permitted. The groupoid layer retains symmetries and comparison witnesses. The stack layer glues coherent local objects on a specified site. The analytic layer asks how large a correction is required and how that correction changes when the model changes. The approximation layer then asks whether the candidates themselves converge. The bounded-budget extension asks whether exactness is possible without unbounded change, and connects a regularized optimizer to a directly checkable budget witness. The phase model and the escaping Hilbert model show why coherent composition, approximate correction, and intended realization remain different questions.

These layers share comparison maps, not an unrestricted identity. Passing from one to another requires the hypotheses printed in the monograph. A coherent family does not automatically realize every intended semantic object, and a theorem about one declared norm is not automatically a theorem about another.

\begin{insight}
\textbf{A practical reading rule}\par
Before applying a Komorebi Space result, identify the views, the restrictions, the protected data, the permitted corrections, and the kind of conclusion required. Existence, uniqueness, symmetry, convergence, and stable recovery have different mathematical tests.
\end{insight}

\subsection*{Companion materials}
The monograph develops the definitions, proofs, worked examples, and background references. The Lean companion maps mathematical topics to their formal declarations and gives reproduction instructions. The accompanying sources include the figures, exact numerical checks, and an executable preservation example. Universal formal theorems, printed arguments for explicit infinite models, and finite computational checks each retain their stated scope.
\end{document}
